\newcommand*{\version}{0.3.0}

\documentclass[12pt,
              article,
              oneside
              ]{memoir}

\usepackage{iftex}

\ifluatex
  \usepackage{fontspec}
\fi

\usepackage[
  % latin,
  english,
  % spanish,% <- principal
  % es-minimal,
  % es-noshorthands,
  % es-tabla,
]{babel}

\ifluatex
% LuaLaTeX --------------------
\defaultfontfeatures{Ligatures=TeX, Scale=MatchLowercase}
\usepackage[MainFeatures={medium,
  Ligatures=Common,
  Numbers=OldStyle,
  SmallCapsFeatures = {LetterSpace=5.0}}]
{junicodevf}
%
% \setsansfont{Libertinus Sans}
% (sin las ligaduras de TeX: -{}-help no se vuelve una raya)
\setmonofont[Scale=MatchLowercase, Ligatures=ResetAll]{DejaVu Sans Mono}
% \setmonofont[Scale=MatchLowercase]{Libertinus Mono}
% \setmonofont[Scale=MatchLowercase}]{Fira Code}
\else
% PDFLaTeX ---------------------
\usepackage[T1]{fontenc}
\usepackage[osf]{libertinus}
\usepackage[scaled=0.9]{FiraMono}
% \usepackage[scaled=0.9]{cascadia-code}
% \usepackage[scaled=0.9]{roboto}
\input glyphtounicode
\pdfgentounicode=1
\fi

\usepackage[babel=true, verbose=false, tracking=true, expansion=true, protrusion=true, final, draft=false]{microtype}
\SetTracking{ encoding = *, shape = sc }{ 30 }
\SetTracking[context = notracking, ]{encoding = *}{0}

\usepackage[autostyle=false, style=british]{csquotes}
\usepackage{latexcolors}

\newcommand{\Rojo}{\color[HTML]{8B0000}}
\newcommand*{\rojo}[1]{\textcolor[HTML]{8B0000}{#1}}
\newcommand*{\rojoit}[1]{\textit{\textcolor[HTML]{8B0000}{#1}}}
\newcommand{\Azul}{\color{bluenivaca}}
\newcommand*{\azul}[1]{\textcolor{bluenivaca}{#1}}
\newcommand{\Verde}{\color{darkgreen}}
\newcommand*{\verde}[1]{\textcolor{darkgreen}{#1}}
\newcommand{\Amarillo}{\color{goldenpoppy}}
\newcommand*{\amarillo}[1]{\textcolor{goldenpoppy}{#1}}
\newcommand{\Morado}{\color{byzantium}}
\newcommand*{\morado}[1]{\textcolor{byzantium}{#1}}
\newcommand{\Naranja}{\color{orange}}
\newcommand*{\naranja}[1]{\textcolor{orange}{#1}}

% \usepackage[margin=3cm]{geometry}

\usepackage{graphicx}
\usepackage{booktabs}

\clubpenalty=9999
\widowpenalty=9999
\brokenpenalty=5000
\frenchspacing
\raggedbottom
\midsloppy
\parindent=14pt
\fboxsep=2pt

\renewcommand\thesection{\arabic{section}}
\chapterstyle{reparticle}

\usepackage{enumitem}
\setlist{noitemsep, nosep}
\setlist[description]{font=\itshape}

\newcommand*{\elipsis}{\,.\,.\,.\,}
\newcommand*{\elipsisb}{\textup{[\elipsis]}}

\renewenvironment{quote}%
 {\begin{list}{}%
              {\setlength\rightmargin{0pt}%
              \setlength\leftmargin{\parindent}}%
  \item[]\small\ignorespaces}
 {\unskip\end{list}}


\definecolor{Link}{HTML}{408080}
\usepackage{hyperref}
\hypersetup{
  colorlinks=true,
  linkcolor=blue,
  filecolor=blue,
  citecolor = black,
  urlcolor=Link,
}
\urlstyle{tt}
\usepackage{bookmark}

\setlength{\droptitle}{-2cm}

\title{\bfseries%
  \texttt{py-tuner} v.\version{} \\
  User Manual}

\author{%
  Nicolás Vaughan\\
  \begin{minipage}[t]{0.8\linewidth}
    \centering\small%
    \href{mailto:nicolasvaughan@proton.me}{\texttt{nicolasvaughan@proton.me}}\\
    \url{https://nicolasvaughan.org}\\
    \url{https://github.com/nivaca/py-tuner}\\
    \url{https://codeberg.org/nivaca/py-tuner}\\
  \end{minipage}}


\date{2026-11-27}

% The pictures, made by tools/manual_images.py
\graphicspath{{img/manual/}{img/}}
\usepackage{longtable}

% the title and the author in the PDF's own metadata
\hypersetup{pdftitle={py-tuner User Manual}, pdfauthor={Nicolás Vaughan}}

% the index (makeindex, with the style manual.ist: see latexmkrc)
\makeindex

% menus, keys, and the pictures of the window
\newcommand*{\menu}[1]{\textit{#1}}
\newcommand*{\then}{\,>\,}
\newcommand*{\key}[1]{\texttt{#1}}
\newcommand*{\shot}[3]{%
  \begin{figure}[htbp]
    \centering
    \includegraphics[width=#2\textwidth,height=0.75\textheight,keepaspectratio]{#1}
    \caption{#3}
    \label{fig:#1}
  \end{figure}}

\begin{document}
\maketitle

\tableofcontents*

\section{Introduction}
\noindent
\texttt{py-tuner} is a tuner for guitar, bass, double bass and violin, with
a chromatic mode for any other instrument. It listens to the sound card,
recognizes the note played, compares it with the nearest string of the
chosen tuning (or with the string one locks) and shows how many cents it
is off. A cent is a hundredth of a semitone.

\index{accuracy}\index{inharmonicity}%
I made it to be accurate. On a steady string, each analysis frame falls
within a few hundredths of a cent of the true pitch, and the reading shown
averages half a second of frames. The pitch is found from all the
harmonics of the string at once, taking into account that a real string is
inharmonic (its upper partials sound a little sharp); so it also works when
a small microphone hardly picks up the fundamental of a low string, as
happens with the bass and the double bass.

\index{interfaces}%
It shows the deviation with a needle and with a strobe, and it runs in
three ways: in a window, in a full-screen terminal and on the command line.
The interface is in English or Spanish. This manual describes version
\version{}; the pictures are of the English window.

\section{Installing}
\label{sec:installing}

\index{installing}%
The releases are on GitHub and Codeberg (the addresses are on the first
page). Each one has the files below, and every file is signed: see
section~\ref{sec:signatures}.

\subsection{Fedora}

\index{fedora@Fedora}\index{copr@Copr}\index{rpm@RPM}%
The RPM is also in Copr, so that \texttt{dnf upgrade} brings each new
version:

\begin{quote}
\texttt{sudo dnf copr enable nivaca/py-tuner}\\
\texttt{sudo dnf install py-tuner}
\end{quote}

\noindent It uses Fedora's own NumPy, PortAudio and PySide6, and installs
a menu entry, the manual pages (\texttt{man py-tuner}, in English and in
Spanish) and this manual.

\subsection{Any Linux: the AppImage}

\index{appimage@AppImage}%
\texttt{py-tuner-\version-x86\_64.AppImage} carries its own Python, Qt and
PortAudio, and needs no installing. It runs on any Linux with glibc 2.34
or later (Ubuntu 22.04, Debian 12, Fedora 35 and later):

\begin{quote}
\texttt{chmod +x py-tuner-\version-x86\_64.AppImage}\\
\texttt{./py-tuner-\version-x86\_64.AppImage}\\
\texttt{./py-tuner-\version-x86\_64.AppImage tui}
\end{quote}

\noindent The first opens the window, the second the terminal interface;
any other command and option of section~\ref{sec:cli} can be given in the
same way.

\subsection{From the source}

\index{source}\index{python@Python}%
With Python 3.14 or later and PortAudio from the system (\texttt{sudo dnf
install portaudio} on Fedora, \texttt{sudo apt install libportaudio2} on
Debian and Ubuntu), the wheel of a release installs with \texttt{pip
install py\_tuner-\version-py3-none-any.whl}. From a copy of the source,
with uv (\url{https://docs.astral.sh/uv/}), \texttt{uv run py-tuner} runs
it at once.

\subsection{Checking the files of a release}
\label{sec:signatures}

\index{signatures}%
Each file of a release has a signature (\texttt{FILE.asc}), and
\texttt{SHA256SUMS} lists the sums of all of them, with its own signature.
They are checked with the public key published with the release:

\begin{quote}
\texttt{gpg -{}-import RPM-GPG-KEY-py-tuner}\\
\texttt{gpg -{}-verify FILE.asc FILE}
\end{quote}

\noindent The RPM is signed inside too; \texttt{dnf} checks it once the
key is imported (\texttt{sudo rpm -{}-import RPM-GPG-KEY-py-tuner}).

\section{Starting the program}

\index{starting}%
From the menu, \menu{py-tuner} opens the window. From a terminal:

\begin{quote}
\begin{tabular}{@{}ll@{}}
\texttt{py-tuner} & the window\\
\texttt{py-tuner tui} & the terminal interface (section~\ref{sec:tui})\\
\texttt{py-tuner cli} & one line of text (section~\ref{sec:cli})\\
\texttt{py-tuner tone E2} & a reference note (section~\ref{sec:tone})\\
\texttt{py-tuner devices} & the list of sound inputs\\
\texttt{py-tuner instruments} & the list of instruments and tunings\\
\end{tabular}
\end{quote}

\noindent The window opens in chromatic mode (section~\ref{sec:chromatic})
until another instrument is chosen; the choice is kept.
\texttt{py-tuner -{}-help} lists every option. The options
change the settings for that session only, unless \texttt{-{}-save} is
added (section~\ref{sec:settings}).

\section{The window}
\label{sec:window}

\index{window}%
\shot{window_en}{0.75}{The window: the G string of a guitar, 0.4 cents sharp}

\noindent From the top down, the window has:

\begin{description}
\item[The bar] the instrument, its tuning (with its strings), the
  reference A (section~\ref{sec:a4}) and the settings.
\item[The note] the note the sound is compared with, in large letters;
  under it, the deviation in cents, and the frequency heard, with an arrow
  to the frequency it should have.
\item[The needle] the deviation, from $-50$ to $+50$ cents. The narrow
  green zone in the middle is the tolerance: a note within it counts as in
  tune (section~\ref{sec:settings}).
\item[The strobe] three bands of stripes that drift to the right when the
  note is sharp and to the left when it is flat, and stand still when it is
  in tune (section~\ref{sec:strobe}).
\item[The history] the deviation over the last ten seconds, the newest at
  the right, over the band of the tolerance (section~\ref{sec:history}).
\item[The strings] a button for each string of the tuning. The string
  compared against is framed.
\item[The status bar] the input in use, its sample rate, and the level of
  the sound, with a mark at the level below which sounds are ignored.
\end{description}

\index{colours}%
The colours say the same everywhere: green within the tolerance, amber
within five times the tolerance, red beyond it.

\shot{flat_en}{0.75}{The low E of a five-string bass, 12.3 cents flat}

\index{string, locked}%
The string compared against is, normally, the nearest one to the sound. A
click on the button of a string locks it: from then on every sound is
compared with that string, however far, which helps when a string is very
slack and sounds nearer to the next one. Another click on the same button
unlocks it.

\shot{locked_en}{0.75}{The violin's A string locked, with A4 at 442~Hz}

\section{Instruments and tunings}
\label{sec:tunings}

\index{instruments}\index{tunings}%
Each instrument has its usual tuning first, and some others:

\begin{longtable}{@{}lll@{}}
\toprule
Instrument & Tuning & Strings, lowest first\\
\midrule
\endhead
Guitar & Standard & \texttt{E2 A2 D3 G3 B3 E4} \\
 & Drop D & \texttt{D2 A2 D3 G3 B3 E4} \\
 & Half step down & \texttt{Eb2 Ab2 Db3 Gb3 Bb3 Eb4} \\
 & Whole step down & \texttt{D2 G2 C3 F3 A3 D4} \\
 & Drop C & \texttt{C2 G2 C3 F3 A3 D4} \\
 & DADGAD & \texttt{D2 A2 D3 G3 A3 D4} \\
 & Open D & \texttt{D2 A2 D3 F\#3 A3 D4} \\
 & Open G & \texttt{D2 G2 D3 G3 B3 D4} \\
 & Open E & \texttt{E2 B2 E3 G\#3 B3 E4} \\
 & Open C & \texttt{C2 G2 C3 G3 C4 E4} \\
 & 7-string & \texttt{B1 E2 A2 D3 G3 B3 E4} \\
\addlinespace
Bass & Standard & \texttt{E1 A1 D2 G2} \\
 & Drop D & \texttt{D1 A1 D2 G2} \\
 & Half step down & \texttt{Eb1 Ab1 Db2 Gb2} \\
 & 5-string & \texttt{B0 E1 A1 D2 G2} \\
 & 5-string high C & \texttt{E1 A1 D2 G2 C3} \\
 & 6-string & \texttt{B0 E1 A1 D2 G2 C3} \\
\addlinespace
Double bass & Orchestral & \texttt{E1 A1 D2 G2} \\
 & Solo & \texttt{F\#1 B1 E2 A2} \\
 & 5-string (low B) & \texttt{B0 E1 A1 D2 G2} \\
 & 5-string (low C) & \texttt{C1 E1 A1 D2 G2} \\
 & Fifths & \texttt{C1 G1 D2 A2} \\
\addlinespace
Violin & Standard & \texttt{G3 D4 A4 E5} \\
 & Cross (AEAE) & \texttt{A3 E4 A4 E5} \\
 & Scordatura GDGD & \texttt{G3 D4 G4 D5} \\
\bottomrule
\end{longtable}

\noindent The instrument also fixes the range in which the fundamental is
looked for, a little beyond its lowest string, so that a very slack string
is still found.

\index{tunings, own}%
\label{sec:own-tunings}
Tunings of one's own are added to the settings file
(section~\ref{sec:files}), with the program closed, as a list
\texttt{custom\_tunings}. Each tuning names its instrument
(\texttt{guitar}, \texttt{bass}, \texttt{double-bass} or
\texttt{violin}), its name and its strings, lowest first:

\begin{verbatim}
"custom_tunings": [
  {"instrument": "guitar", "name": "Nashville",
   "strings": "E3 A3 D4 G4 B3 E4"}
]
\end{verbatim}

\noindent They appear after the tunings of the instrument. A tuning with a
note that cannot be read is left out.

\section{The reference A}
\label{sec:a4}

\index{reference A}\index{a4@A4}%
All the notes are measured against the A above middle C (A4), 440~Hz
unless changed. It can be set from 380 to 480~Hz, in tenths of a hertz: in
the window, with the box \menu{A4} (the arrow keys move it by 0.1~Hz, Page
Up and Page Down by 1~Hz) or with the menu beside it, which has the usual
ones (415, 430, 432, 435, 438 to 444, 446 and 466~Hz); in the terminal, with
the arrow keys; on the command line, with \texttt{-a}:

\begin{quote}
\texttt{py-tuner -a 442}
\end{quote}

\noindent Raising A4 by 1~Hz raises every note by about 3.9 cents.

\section{Chromatic mode and note names}
\label{sec:chromatic}

\index{chromatic mode}%
The instrument \menu{Chromatic}, the one the program starts with, has no
strings: the sound is compared with the nearest note of the equal-tempered
scale, from A0 to G6. It serves for
any instrument, and for checking the notes along a string.

\index{note names}%
The notes can be written in English (C, D, E\elipsis), Latin (Do, Re,
Mi\elipsis) or German (with H for B and B for B$\flat$), with sharps,
flats or both (F$\sharp$/G$\flat$); the octave is always the scientific
one, in which middle C is C4.

\shot{chromatic_en}{0.75}{Chromatic mode, with Latin names and flats}

\section{The needle and the strobe}
\label{sec:strobe}

\index{needle}%
The needle shows the deviation at once, and is the easier to read from
afar. It follows the reading with a little inertia, so that it does not
tremble.

\index{needle!fine scale}%
Its scale is linear unless changed: a cent is a cent anywhere on it. On the
fine scale, chosen in the settings, with \key{g} in the terminal or with
\texttt{-{}-needle fine}, the last five cents either side take half of each
side, and the rest up to fifty the other half; the last cents of the tuning
can then be told apart, while a string far off is still on the scale. The
idea is Tuneo's (Braulio Ríos), whose gauge is drawn so
(figure~\ref{fig:spectrum_en}).

\index{strobe}%
The strobe is the more sensitive, and it is a real one. A mechanical strobe
turns at the frequency of the note wanted and is lit by the sound: in tune,
its pattern stands still. Each band of py-tuner's does the same with
numbers: it follows the phase of one partial of the string, and moves one
stripe per second for each cent the partial is off, to the right when it is
sharp, to the left when it is flat. So a deviation of a tenth of a cent,
which the needle hardly shows, still moves a band slowly. A band fades when
its partial is lost in the noise. Its idea comes from the stroboscope of
Entropy Piano Tuner (Haye Hinrichsen and Christoph Wick) and from Strobie
(Davorin Šego).

\index{strobe!vernier}\index{strobe!harmonics}%
It has two modes, chosen in the settings or with \key{m} in the terminal:

\begin{description}
\item[Vernier] (the usual) the three bands follow the clearest partial, the
  top one at its own speed, the middle one twice as fast and the bottom one
  four times. To tune, turn the peg until the bottom band stands still.
\item[Harmonics] each band follows a partial of its own: the fundamental,
  the second and the third. A real string is a little stiff, and its upper
  partials run sharp; py-tuner looks for each where the string's own
  stiffness puts it, so on a string in tune the three stand still.
\end{description}

\noindent Which of the needle and the strobe is shown is chosen in the
settings, or with \key{v} in the terminal.

\shot{strobe_en}{0.75}{The strobe alone}

\section{The history}
\label{sec:history}

\index{history}\index{trace}%
Under the strobe, the history shows the deviation over the last ten
seconds, the newest at the right, over the green band of the tolerance: a
string settling after a turn of the peg, or a peg turned the wrong way, is
seen at once. A gap is a moment without a note. It is hidden in the
settings, with \key{h} in the terminal, or with \texttt{-{}-no-trace}.

\section{The spectrum}
\label{sec:spectrum}

\index{spectrum}%
Below the history, the window can show the spectrum of the sound: its
frequencies from 30~Hz to 5~kHz on a logarithmic axis (the lines are the
A's), its clearest peaks named as notes with their cents, and the harmonics
of the note sought marked in green. It shows what the tuner hears: a weak
fundamental, a buzz, another string ringing, the hum of the room. In the
terminal it is a row of bars with the peaks named under it. It is shown in
the settings, with \key{f} in the terminal or with \texttt{-{}-spectrum}.
Strobie (Davorin Šego) has one so.

\shot{spectrum_en}{0.75}{The fine scale of the needle and the spectrum: an A2 7.5 cents flat}

\section{Offsets, transposition and capo}
\label{sec:offsets}

\index{offsets}%
Some players tune a few notes a little off equal temperament, so that the
chords they play most sound sweeter. \menu{Settings\then Note offsets\elipsis}
gives each note name its offset, from $-50$ to $+50$ cents, in every octave;
the needle, the strobe and the tone then aim at the note so moved. On the
command line: \texttt{-{}-offset B=-2} (repeated for each note;
\texttt{B=0} removes it).

\shot{offsets_en}{0.35}{Note offsets: B 2 cents flat, G 1.5 sharp, E 1 flat}

\index{transposition}%
\menu{Transpose} shows the names moved some semitones, for a transposing
instrument: 2 for one in B$\flat$, whose written C sounds B$\flat$. Only the
names change; the pitch measured is the same. \index{capo}%
\menu{Capo} is the fret of a capo: every string is raised as many
semitones, and the buttons of the strings show the notes they then sound.
On the command line, \texttt{-{}-transpose} and \texttt{-{}-capo}; in the
terminal, \key{c} moves the capo.

\section{The reference tone}
\label{sec:tone}

\index{reference tone}\index{tone}%
To tune by ear, py-tuner can play the note through the computer's sound, as
pianito (Axel Stevens) does: a plucked string of a few seconds, at the
exact frequency aimed at, with the reference A, the offset of the note and
the capo. \menu{Tone}, in the window, plays the locked string, or else the
note shown, or else A4; \key{p} does the same in the terminal. On the
command line:

\begin{quote}
\texttt{py-tuner tone E2}\\
\texttt{py-tuner tone -i guitar 1}\quad(the first string)\\
\texttt{py-tuner tone -a 442}\quad(A4 at 442~Hz)
\end{quote}

\section{Settings}
\label{sec:settings}

\index{settings}%
\menu{Settings\elipsis} opens the less frequent settings:

\shot{settings_en}{0.6}{The settings}

\begin{description}
\item[Input device] the sound input to listen to, or the system's default.
\item[Channel] the mix of all channels, or only one (for an interface with
  the instrument in one input).
\item[Note names] English, Latin or German.
\item[Accidentals] sharps, flats or both (F$\sharp$/G$\flat$).
\item[Show] the needle and the strobe, or one of them.
\item[Language] the language of the interface, or the system's; it is used
  from the next time the program starts.
\item[In tune within] the tolerance, in cents (1 unless changed).
\item[Ignore below] sounds quieter than this, in dBFS, are ignored ($-60$
  unless changed). Raise it in a noisy room.
\item[Smoothing] how many seconds the reading averages (0.5 unless
  changed): longer is steadier, shorter follows faster.
\item[Analysis frame] the length of each frame, in periods of the note (16
  unless changed): more is more precise, fewer reacts faster.
\item[Response] normal (unless changed) or steady, for the double bass
  (section~\ref{sec:how}).
\item[Strobe] vernier or harmonics (section~\ref{sec:strobe}).
\item[Needle scale] linear or fine (section~\ref{sec:strobe}).
\item[Show the spectrum] the spectrum and its peaks
  (section~\ref{sec:spectrum}).
\item[Learn the room's noise] (on unless changed) ignore what does not
  stand 10~dB over the background noise, learned while nothing plays
  (section~\ref{sec:how}).
\item[Show the history] the history of the deviation
  (section~\ref{sec:history}).
\item[Transpose, Capo, Offsets] section~\ref{sec:offsets}.
\end{description}

\noindent Every change is saved at once, and the window, the terminal and
the command line share the same settings.

\section{The terminal}
\label{sec:tui}

\index{terminal}\index{tui@TUI}%
\texttt{py-tuner tui} shows, in the whole terminal, the same as the window:
the note in large letters, the frequency and the deviation, a needle of
text, the strobe, the strings, and the level.

\shot{tui_en}{0.9}{The terminal interface, with the G string of a guitar 0.4 cents sharp}

\noindent It is driven with these keys, and each change is saved in the
settings:

\begin{longtable}{@{}lp{0.7\textwidth}@{}}
\toprule
Key & What it does\\
\midrule
\endhead
\key{q} & Quit \\
\key{← / →} & Lower / raise the reference A4 by 0.1 Hz \\
\key{↓ / ↑} & Lower / raise the reference A4 by 1 Hz \\
\key{0} & Put the reference A4 back to 440 Hz \\
\key{i} & The next instrument \\
\key{t} & The next tuning of the instrument \\
\key{s} & Lock the next string (after the last one, back to the nearest string) \\
\key{n} & The next system of note names: English, Latin, German \\
\key{b} & Sharps, flats or both \\
\key{v} & Needle and strobe, the needle alone or the strobe alone \\
\key{m} & The strobe: vernier (one partial, three speeds) or harmonics (partials 1, 2, 3) \\
\key{h} & Show or hide the history of the deviation \\
\key{c} & Move the capo one fret up (after the 7th, off) \\
\key{f} & Show or hide the spectrum, with its peaks named \\
\key{g} & The needle's scale: linear, or fine (the centre spread out) \\
\key{r} & The response: normal, or steady (for the double bass) \\
\key{p} & Play the reference tone: the locked string, or the note shown, or A4 \\
\key{l} & The next language of the interface \\
\bottomrule
\end{longtable}

\section{The command line}
\label{sec:cli}

\index{command line}\index{cli@CLI}%
\texttt{py-tuner cli} writes the reading in a line that rewrites itself
(or, in a pipe, a line every quarter of a second):

\begin{verbatim}
A2       109.52 Hz   -7.5 cents  [---------*|----------]  string 2
\end{verbatim}

\noindent \texttt{-{}-once} waits for the first steady reading, writes it
and ends, for use in scripts.

\shot{cli_en}{0.95}{The command line: a recording measured note by note, and the first steady reading}

\index{recordings}\index{wav@WAV}%
With a WAV file, \texttt{py-tuner cli -{}-file FILE.wav} analyses the whole
recording at once and writes a line for each note heard in it: when it
sounds, the note, its frequency, the frequency it should have, the
deviation and its spread:

{\footnotesize
\begin{verbatim}
0.20-   1.77 s  E2  82.550 Hz  (target 82.407)  +3.00 cents  ±0.02  string 1
1.80-   3.52 s  A2  109.52 Hz  (target 110.00)  -7.50 cents  ±0.03  string 2
\end{verbatim}
}

\noindent The first tenth of a second of each note is left out, since a
plucked string sounds a little sharp at the start. With
\texttt{-{}-realtime}, the file is played at its own pace instead, as if
it came from the sound card; \texttt{-{}-file} works in the same way with
the window and the terminal.

The options common to the three ways:

\begin{longtable}{@{}lp{0.62\textwidth}@{}}
\toprule
Option & What it does\\
\midrule
\endhead
\texttt{-a}, \texttt{-{}-a4 HZ} & the reference A4, 380 to 480 Hz\\
\texttt{-i}, \texttt{-{}-instrument} & \texttt{chromatic} (unless changed), \texttt{guitar}, \texttt{bass}, \texttt{double-bass} or \texttt{violin}\\
\texttt{-t}, \texttt{-{}-tuning NAME} & the tuning, by its name (e.g.\ \texttt{"Drop D"})\\
\texttt{-s}, \texttt{-{}-string N} & compare only with string N (1 = the lowest); \texttt{auto}, with the nearest\\
\texttt{-{}-names} & \texttt{english}, \texttt{latin} or \texttt{german}\\
\texttt{-{}-accidentals} & \texttt{sharps}, \texttt{flats} or \texttt{both} (F\#/Gb); also \texttt{-{}-flats}, \texttt{-{}-sharps}\\
\texttt{-{}-tolerance CENTS} & in tune within this many cents\\
\texttt{-d}, \texttt{-{}-device DEV} & the input, by its number (\texttt{py-tuner devices}) or part of its name\\
\texttt{-r}, \texttt{-{}-rate HZ} & the sample rate (the device's, unless given)\\
\texttt{-c}, \texttt{-{}-channel} & \texttt{mix}, or a channel number from 1\\
\texttt{-f}, \texttt{-{}-file WAV} & a WAV file instead of the sound card\\
\texttt{-{}-min-level DBFS} & ignore sounds quieter than this\\
\texttt{-{}-smoothing SECONDS} & seconds the reading averages\\
\texttt{-{}-periods N} & periods of the note in each analysis frame\\
\texttt{-{}-response} & \texttt{normal} or \texttt{steady} (for the double bass)\\
\texttt{-{}-display} & \texttt{both}, \texttt{needle} or \texttt{strobe}\\
\texttt{-{}-offset NOTE=CENTS} & tune a note name off equal temperament, in every octave\\
\texttt{-{}-transpose N} & the names for a transposing instrument\\
\texttt{-{}-capo FRET} & a capo, which raises every string\\
\texttt{-{}-strobe} & \texttt{vernier} or \texttt{harmonics}\\
\texttt{-{}-trace}, \texttt{-{}-no-trace} & the history of the deviation\\
\texttt{-{}-spectrum} & the spectrum, with its peaks named (\texttt{-{}-no-spectrum} hides it)\\
\texttt{-{}-needle} & \texttt{linear} or \texttt{fine}\\
\texttt{-{}-learn-noise} & ignore what does not stand over the room's noise\\
\texttt{-{}-no-learn-noise} & only the fixed level of \texttt{-{}-min-level}\\
\texttt{-{}-language} & \texttt{en} or \texttt{es}\\
\texttt{-{}-save} & keep the options given as the new settings\\
\bottomrule
\end{longtable}

\section{How it measures}
\label{sec:how}

\index{accuracy}%
Each analysis takes two steps. A coarse one finds the period of the sound
by McLeod's method, which cannot be wrong by an octave; a fine one then
locates every harmonic of the string, up to the sixteenth, at the exact
peak of its spectrum, and fits all of them at once to the model of a real,
slightly stiff string, $f_n = n f_0 \sqrt{1 + B n^2}$. Strong harmonics
count more than weak ones, and those buried in noise are left out. The
reading shown is the trimmed mean of the last half second of frames.

\index{noise}%
Before any of this, a sound must be loud enough. Besides the fixed level of
\menu{Ignore below}, py-tuner learns the noise of the room while nothing
plays, as Strobie does: the floor follows the level when nothing sounds,
only creeps up under a sound, and never rises during a note, which would
otherwise become the background. A sound must stand 10~dB over that floor.
The level bar marks the threshold in use.

\index{octave}%
A real string can fool a period: two strings ringing in sympathy share a
period an octave or a twelfth below the note, and a steel string whose
fundamental the microphone hardly picks up repeats almost every half
period. So the candidates around the period (a third, a half, twice and
three times it) are scored by how far their harmonics stand over the noise
around them, less what stands where a wrong octave's harmonics would fall,
an idea taken from Strobie and from Entropy Piano Tuner; the note already
sounding is favoured while it decays, as Tuneo (Braulio Ríos) does. A new
note right after a pluck is taken at once; without a pluck, it must last
0.3 seconds, or 0.6 if it is an octave, a twelfth or a fifth of the note
sounding. And with an instrument's strings, only the notes near them are
looked for. On eleven recordings of real guitar strings, these rules left
one false note where there were fifteen.

\index{double bass}\index{response}%
The pitch of a low string is not still. A plucked double bass string is
stretched by its own swing, so it starts sharp, several cents if it is
plucked hard, and goes down as it fades; the bow moves it a little too.
The normal response shows all of this, and the needle wanders with it.
The steady response (\menu{Response} in the settings, \key{r} in the
terminal, \texttt{-{}-response steady}) leaves out the first half second
of each pluck, while the note that was sounding stays on the screen;
averages at least a second and a half (more, if \menu{Smoothing} says
more); moves the needle more slowly; and, with an instrument's strings,
keeps reading the faint tail of the note while the other strings ring
with it. A note that rises from silence without a pluck, a bowed one for
instance, then waits 0.3~s before it is shown. On synthetic double bass
strings, plucked and bowed, the reading moved a fifth to a third as much.

\index{tips}%
A few habits help:

\begin{itemize}
\item Play the string alone, and let the first instant of the pluck pass:
  a plucked string sounds a little sharp at first, and settles.
\item Bring the microphone near the instrument, away from fans and other
  noise; if the room is noisy, raise \menu{Ignore below}.
\item For low strings, give the reading half a second more: the frame of a
  low E of a bass lasts almost 0.4~s. On a double bass, use the steady
  response.
\item To tune finely, watch the bottom band of the strobe.
\end{itemize}

\section{The files it keeps}
\label{sec:files}

\index{files}%
The program keeps one file: \texttt{\textasciitilde/.config/py-tuner/settings.json}
(under \texttt{\$XDG\_CONFIG\_HOME}, if it is set), with every setting, in
JSON. It can be edited by hand with the program closed; a value that does
not make sense is replaced by its default, so the program always starts.
Deleting it brings back all the defaults.

\section{When something goes wrong}

\index{problems}%
\begin{description}
\item[Nothing is heard (``Play a note'')] the level bar does not move: the
  input chosen is not the one with the instrument. Choose another in
  \menu{Settings} (\texttt{py-tuner devices} lists them), and check that
  the input is not muted in the system's sound settings.
\item[The level moves, but no note appears] the sound is quieter than
  \menu{Ignore below}, or too noisy to have a pitch. Bring the microphone
  nearer, or lower \menu{Ignore below}.
\item[The note jumps to the next string] the string is so slack that it
  is nearer to another; lock the string you are tuning
  (section~\ref{sec:window}).
\item[The reading is a few cents off from another tuner] check that both
  have the same reference A.
\item[``Cannot open the input''] another program holds the sound card, or
  the device was unplugged; the message gives the reason.
\item[The faint tail of a string ends too soon] the tail has sunk into the
  noise of the room; turn off \menu{Learn the room's noise} to read it to
  the end.
\item[The tone does not sound] the computer's sound output is muted, or
  another program holds it; the message says why.
\item[A setting is wrong] delete the settings file
  (section~\ref{sec:files}).
\end{description}

\noindent
Problems and wishes can be written as issues on GitHub or
Codeberg (the addresses are on the first page), or sent to
\href{mailto:nicolasvaughan@proton.me}{\texttt{nicolasvaughan@proton.me}}.

\printindex

\end{document}

% ===============================================
%%% Local Variables:
%%% mode: latex
%%% TeX-master: t
%%% End:
